\documentclass[book.tex]{subfiles}

\begin{document}

\chapter{Poisson Brackets and Symplectic Manifolds}
\label{chap:poisson_bracket}
\noindent\rule{11cm}{0.4pt} \\
\textbf{Prerequisites:} \autoref{chap:manifold} \\
\textbf{Difficulty Level:} *** \\
\noindent\rule{11cm}{0.4pt}
\vspace{5mm}

Symplectic manifolds arise in classical mechanics as a useful model of the phase space of the system.

\section{Mathematical Definitions}
A symplectic manifold is a smooth manifold $M$ that has a closed non-degenerate differential 2-form $\omega$. Formally, a 2-form is a function from the manifold to the second exterior power of the cotangent space for the manifold.
\begin{align}
\omega: M \to \wedge^2 T_p^* M
\end{align}


Non-degenerate means that if there exists $X \in T_p M$ such that $\omega(X, Y) = 0$ for all $Y \in T_p M$ then $X = 0$.

%\textcolor{red}{TODO: Add an explanation of differential forms to the manifolds chapter.}

Let $M$ be a smooth manifold of dimension $n$. Consider the cotangent bundle $T^*M$. Then the canonical symplectic form is given by
\begin{align}
\omega = \sum_{i=1}^n dp_i \wedge dq^i
\end{align}
Here $(q^1,\dotsc,q^n)$ are local coordinates on $M$ and $(p_1,\dotsc,p_n)$ are fiberwise coordinates with respect to cotangent vectors $dq^1,\dotsc,dq^n$.

%\textcolor{red}{TODO: This explanation needs to be fleshed out.}

%\section{Darboux Coordinates}

%TODO: Explain this. Darboux coordinates are the "right" way of getting the canonical $p$, $q$ coordinates

\section{Poisson Bracket}
A Poisson Bracket is a function, typically denoted by $\{\cdot, \cdot \}$ that satisfies the following properties.

\begin{itemize}
\item Anticommutativity: 
\begin{align}
\{f, f\} &= - \{g, f\}
\end{align}
\item Bilinearity: For $a, b \in \mathbb{R}$
\begin{align}
\{af + bg, h\} &= a\{f, h\} + b \{g, h\}\\
\{h, af + bg\} &= a\{h, f\} + b \{h, g\}
\end{align}
\item Leibniz's Rule
\begin{align}
\{fg, h\} &= \{f, h\}g + f\{g, h\} 
\end{align}
\item Jacobi Identity:
\begin{align}
\{f, \{g, h\}\} + \{g, \{h, f\}\} + \{h, \{f, g\}\} &= 0
\end{align}
\end{itemize}

%In Darboux coordinates, the Poisson bracket 
In canonical coordinates, the Poisson bracket takes the form
\begin{align}
\{f, g\} &= \sum_{i=1}^N \left ( \frac{\partial f}{\partial q_i}\frac{\partial g}{\partial p_i} - \frac{\partial f}{\partial p_i} \frac{\partial g}{\partial q_i} \right )
\end{align}
%\section{Exercises}
%\begin{enumerate}
%    \item TODO
%\end{enumerate}

\end{document}